\renewcommand{\theequation}{A\thechapter.\arabic{equation}}
\label{ref:appendixA}

This derivation begins with the Equations \eqref{eq:wilson-R} and \eqref{eq:wilson-n} provided by Wilson and England \cite{wilson2002delayed}:
\begin{equation}
    R(t) = \sum_{k=1}^{K} \bar{\nu}_{d, k} \lambda_k e^{-\lambda_k t},
    \label{eq:wilson-R}
\end{equation}
and,
\begin{equation}
    \dot{n}_d(T, t) = \bar{\nu}_d \int_{0}^{T} F(\tau) R(T-\tau+t) d\tau,
    \label{eq:wilson-n}
\end{equation}
where,
\begin{align}
T &= \mbox{ the length of the fission history [$s$]} \nonumber \\
t &= \mbox{ the decay time from the end of the fission history [$s$]} \nonumber \\
F(\tau) &= \mbox{ the fission rate history [$\frac{fissions}{s}$].} \nonumber
\end{align}
Because the results are normalized to a single delayed neutron, the $\bar{\nu}_d$ term is divided out.

To derive an analytical solution to this problem for the 0D flow approach where the in-core and ex-core residence times, $\tau_{in}$ and $\tau_{ex}$, are equivalent, an approximation of the fission history must be made.
For this derivation, $\tau_{res}$ will be used to represent both the in-core and ex-core residence times.
Because the 0D flow model's fission history is a positive square wave, analytical integration is not feasible.
Instead, a Fourier decomposition of the square wave into a sum of sine waves makes the problem far more approachable.
This is given as:
%For numerical integration this is not an issue, but an analytical solution is useful for setting up the least squares matrices.
\begin{equation}
    \frac{F(t)}{F_0} = \frac{1}{2} + \frac{2}{\pi} \left(sin \left(\frac{\pi t}{\tau_{res}} \right) + \frac{1}{3} sin \left(\frac{3 \pi t}{\tau_{res}} \right) + \dots  \right),
    \label{eq:pos-square}
\end{equation}
where,
\begin{align}
    F_0 =& \text{ the magnitude of the constant fission rate [$\frac{fissions}{s}$].} \nonumber
\end{align}

Plugging this in and separating the integral, for group $k$ and for $j$ Fourier terms, gives:
\begin{equation}
\begin{split}
    \dot{n}_{d, k}(T, t) =& \; \frac{F_0 \bar{\nu}_{d, k} \lambda_k}{2} \int_{0}^{T} e^{-\lambda_k \left( T-\tau + t \right)} d\tau\\
    & + \frac{2 F_0 \bar{\nu}_{d, k} \lambda_k}{\pi} \int_{0}^{T} sin \left( \frac{\pi \tau}{\tau_{res}}  \right) e^{-\lambda_k \left( T-\tau + t \right)} d\tau\\
    & + \dots\\
    & + \frac{2 F_0 \bar{\nu}_{d, k} \lambda_k}{(2j+1) \pi} \int_{0}^{T} sin \left( \frac{(2j+1) \pi \tau}{\tau_{res}}  \right) e^{-\lambda_k \left( T-\tau + t \right)} d\tau
    \label{eq:split-integral}
\end{split}
\end{equation}

Solving the integrals and simplifying yields: 
\begin{equation}
\begin{split}
    \dot{n}_{d, k}(T, t) =& \; \frac{F_0 \bar{\nu}_{d, k} e^{-\lambda_k t} \left( 1 - e^{-\lambda_k T} \right)}{2}\\
    & + \dots\\
    & + \frac{2 F_0 \bar{\nu}_{d, k} \lambda_k^2 \tau_{res}^2 e^{-\lambda_k t} }{\left( (2j+1) \pi \right) \left( (2 \pi j + \pi)^2 + \lambda_k^2 \tau_{res}^2 \right)} sin\left( \frac{\pi (2j+1) T}{\tau_{res}} \right)\\
    & + \frac{2 F_0 \bar{\nu}_{d, k} \lambda_k \tau_{res} e^{-\lambda_k t} e^{-\lambda_k T}}{\left( (2 \pi j + \pi)^2 + \lambda_k^2 \tau_{res}^2 \right)} cos\left( \frac{\pi (2j+1) T}{\tau_{res}} \right)  \\
    & + \frac{2 F_0 \bar{\nu}_{d, k} \lambda_k \tau_{res} e^{-\lambda_k t} e^{-\lambda_k T}}{\left( (2j+1) \pi \right) \left( (2 \pi j + \pi)^2 + \lambda_k^2 \tau_{res}^2 \right)} 2 \pi j  \\
    & + \frac{2 F_0 \bar{\nu}_{d, k} \lambda_k \tau_{res} e^{-\lambda_k t} e^{-\lambda_k T}}{\left( (2j+1) \pi \right) \left( (2 \pi j + \pi)^2 + \lambda_k^2 \tau_{res}^2 \right)} \pi,
    \label{eq:split-integral-solved}
\end{split}
\end{equation}
where each integrated Fourier term generates four terms added to the sum.

Similar to \cite{wilson2002delayed}, the irradiation time is taken to the limit as it approaches infinity, though only for the non-periodic terms.
By irradiating for several half-lives of time for each \gls{DNP}, where the longest half-life is roughly one minute, then the irradiation time can be taken as infinite. 
There are two reasons this limit is taken.
Firstly, a pulse irradiation with circulation will not provide any unique results when compared to traditional data which already exists.
Additionally, this allows the equation to be simplified, as shown in Equation \eqref{eq:split-integral-solved-reduced}.
Only the sine term does not go to zero as the irradiation time goes to infinity, so it is the only term significant for saturation irradiation in the 0D flow model.


\begin{equation}
\begin{split}
    \dot{n}_{d, k}(T, t) =& \; \frac{F_0 \bar{\nu}_{d, k} e^{-\lambda_k t}}{2}\\
    & + \dots\\
    & + \frac{2 F_0 \bar{\nu}_{d, k} \lambda_i^2 \tau_{res}^2 e^{-\lambda_k t} }{\left( (2j+1) \pi \right) \left( (2 \pi j + \pi)^2 + \lambda_k^2 \tau_{res}^2 \right)} sin\left( \frac{\pi (2j+1) T}{\tau_{res}} \right)\\
    \label{eq:split-integral-solved-reduced}
\end{split}
\end{equation}

The final assumption to make is that the irradiation time, $T$, is an integer multiple of the residence times, $\tau_{res}$.
This assumption makes the sine term 0 for every value of $k$, which simplifies the result to:
\begin{equation}
    \dot{n}_{d}(t) = \frac{F_0}{2} \sum_{k=1}^{K} \bar{\nu}_{d, k} e^{-\lambda_k t}
    \label{eq:fully-resolved-flow-irrad}
\end{equation}
This is very similar to the saturation irradiation delayed neutron count rate equation, which is:
\begin{equation}
    \dot{n}_{d}(t) = F_0 \sum_{k=1}^{K} \bar{\nu}_{d, k} e^{-\lambda_k t}
    \label{eq:saturation-count-eqn}
\end{equation}

Because of the initial assumption made, which was that half the time is in-core and half ex-core, the average fission rate is half of the fission rate scaling term, $F_0$.
For both cases (constant irradiation and half in-core and half ex-core), the fission term in front is simply the average fission rate over the irradiation time.
The generic case is addressed in Appendix \ref{ref:appendixB}.

\begin{comment}

An additional challenge is the fact that the square wave needs to vary from 0 to 1, rather than -1 to 1.
However, this can be accounted for by making the substitution shown in Equation \eqref{eq:pos_sin}.
By replacing the sine functions with this form, the values will either be positive, or 0 where they would have been negative.

\begin{equation}
    sin_+(x) = \frac{sin(x) + |sin(x)|}{2}
    \label{eq:pos_sin}
\end{equation}

%The final hurdle for this derivation is the variation in the in-core and ex-core residence times.
%This means the fission history shape function needs to vary the sine period based on if it is in-core or ex-core.
%One way to account for this is to include a Heaviside function that takes in the sine function as an argument.
%This is shown in Equation \eqref{eq:first_term}, which shows the first term in the Fourier composition of the square wave.

%\begin{equation}
%    \psi(t) = sin_+\left(\frac{\pi}{\tau} t \right) H\left(sin_+\left(\frac{\pi}{\tau} t\right) \right)
%    \label{eq:first_term}
%\end{equation}

Another issue arises when constructing the square wave, which is that non-zero terms begin appearing in the ex-core residence time sections.
One way to account for this is to include a Heaviside function that takes in the sine function as an argument.
This is shown in Equation \eqref{eq:first_term}, which shows the first two terms in the Fourier composition of the square wave.

\begin{equation}
\begin{split}
    \psi(t) = & \; sin_+\left(\frac{\pi t}{\tau_{res}} \right) H\left[sin\left(\frac{\pi t}{\tau_{res}} \right) \right] + \\
    & \; \frac{1}{3} sin_+\left(\frac{\pi t}{3 \tau_{res}} \right) H\left[sin\left(\frac{\pi t}{\tau_{res}} \right) \right]
    \label{eq:first_term}
\end{split}
\end{equation}

The Heaviside function can be pulled out and moved to the front, yielding Equation \eqref{eq:simplified_term}:

\begin{equation}
    \psi(t) = H\left[sin\left(\frac{\pi t} {\tau_{res}} \right) \right] \left( sin_+\left(\frac{\pi t}{\tau_{res}} \right) + \frac{1}{3} sin_+\left(\frac{\pi t}{3 \tau_{res}} \right) + \dots \right)
    \label{eq:simplified_term}
\end{equation}

This can be plugged back into the fission history term, which can then be used in Equation \eqref{eq:wilson-n}.
To start, group $i$ is considered and only the first Fourier series term is considered.

\begin{equation}
    n_d(T, t) = F_0 \int_{0}^{T} H\left[sin\left(\frac{\pi \tau} {\tau_{res}} \right) \right] \left( sin_+\left(\frac{\pi \tau}{\tau_{res}} \right) \right) a_i \lambda_i e^{-\lambda_i (T-\tau +t)} d\tau
    \label{eq:wilson-n}
\end{equation}

Based on the definition of $sin_+(x)$, the integral can be split into two separate integrals for the absolute and non-absolute forms of sine.
The two integrals are given in Equations \eqref{eq:non-abs} and \eqref{eq:abs}.

\begin{equation}
    \frac{1}{2} \int_0^T  H\left[sin\left(\frac{\pi \tau} {\tau_{res}} \right) \right] sin\left(\frac{\pi \tau} {\tau_{res}}\right) a_i \lambda_i e^{-\lambda_i (T-\tau +t)} d\tau
    \label{eq:non-abs}
\end{equation}

\begin{equation}
    \frac{1}{2} \int_0^T  H\left[sin\left(\frac{\pi \tau} {\tau_{res}} \right) \right] \left|sin\left(\frac{\pi \tau} {\tau_{res}}\right)\right| a_i \lambda_i e^{-\lambda_i (T-\tau +t)} d\tau
    \label{eq:abs}
\end{equation}

To solve these equations, only the regions in which the Heaviside function is non-zero are considered.
This provides bounds, where $\tau$ is between $2 \tau_{res} k$ and $\tau_{res} (2 k + 1) $.
The value of $k$ is an integer which starts at 0 and increases until the irradiation time, $T$, is met by $\tau$.
This occurs at the value of $k$ shown in Equation \eqref{eq:k-time-done}, assuming that each full irradiation period occurs an integer number of times.

\begin{equation}
    k = \frac{T - 2 \tau_{res}}{2 \tau_{res}}
    \label{eq:k-time-done}
\end{equation}






\hline







This means that a single period can be evaluated, and then all that is required is a scaling factor based on the total irradiation time of the sample relative to the period.
Evaluating over the period yields a positive portion of the sine profile followed by a zero portion due to the Heaviside function.
This means that the absolute value of sine will yield the same result as the sine, meaning only a single integral needs to be solved.
This can be seen in Equation \eqref{eq:reduced-eqn}.

\begin{equation}
    \int_0^{\tau_{res}} sin\left(\frac{\pi \tau} {\tau_{res}}\right) a_i \lambda_i e^{-\lambda_i (T-\tau +t)} d\tau
    \label{eq:reduced-eqn}
\end{equation}

Solving this equation yields Equation \eqref{eq:first-term-full}, which shows the full form of the delayed neutron count rate for group $i$ as a function of irradiation time and time after irradiation for the first Fourier component of the square wave.

\begin{equation}
    n_{d, 1}(T, t) = F_0 a_i \lambda_i \left( \frac{\pi \tau_{res} (e^{-\lambda_i \tau_{res}} + 1) e^{-\lambda_i T} e^{-\lambda_i t}}{\lambda_i^2 \tau_{res}^2 + \pi^2} \right)
    \label{eq:first-term-full}
\end{equation}

Similar to \cite{wilson2002delayed}, the irradiation time is taken to the limit as it approaches infinity.
There are three reasons for this.
Firstly, a pulse irradiation with circulation will not provide any unique results when compared to traditional data which already exists.
Secondly, by irradiating for several half-lives of time for each \gls{DNP}, roughly 5 minutes, then the irradiation time can be taken as infinite.
Finally, this allows the equation to be simplified even further, though only slightly.

\begin{equation}
    n_{d, 1, i}(T, t) = F_0 a_i \lambda_i \left( \frac{\pi \tau_{res} (e^{-\lambda_i \tau_{res}} + 1)}{\lambda_i^2 \tau_{res}^2 + \pi^2} \right) e^{-\lambda_i t}
    \label{eq:first-term-full-final}
\end{equation}

To simplify, this term due to the fission history square wave shape for group $i$ for the first Fourier component will be captured by the $\eta_{1, i}$ term, as shown in Equation \eqref{eq:eta-first}.

\begin{equation}
    \eta_{1, i} = \left( \frac{\pi \tau_{res} (e^{-\lambda_i \tau_{res}} + 1)}{\lambda_i^2 \tau_{res}^2 + \pi^2} \right)
    \label{eq:eta-first}
\end{equation}

Equation \eqref{eq:kth-term-full} shows how to calculate the $\eta_{k, i}$ where $k$ is the $k^{th}$ term of the square wave.

\begin{equation}
    \eta_{k, i} = \left( \frac{k \tau_{res} e^{-\lambda_i t} \left( e^{\lambda_i \tau_{res}} \left( k \lambda_i \tau_{res} sin\left(  \frac{\pi}{k}  \right) - \pi cos \left( \frac{\pi}{k}   \right)    \right) + \pi \right)}{k^2 \lambda_i^2 \tau_{res}^2 + \pi^2} \right)
    \label{eq:kth-term-full}
\end{equation}

Because the Fourier series terms are summed, the integrals can be summed separately.
This yields the general equation, shown in Equation \eqref{eq:full-equation-delnugroup} for the delayed neutron count group fit, where $\kappa$ is a number of sine waves such that the square wave is sufficiently well represented.

\begin{equation}
    n_d(t) = F_0 \sum^{I}_{i=1} \left( a_i \lambda_i e^{-\lambda_i t} \sum^{\kappa}_{k=1} \eta_{k, i} \right)
    \label{eq:full-equation-delnugroup}
\end{equation}

\end{comment}